\documentclass{article}
\usepackage{amsmath, amssymb, amsthm}
\usepackage{hyperref}
\usepackage{geometry}
\geometry{margin=1in}

\title{Erd\H{o}s Problem \#413}
\date{Last updated: 2025-08-31}
\author{Source: erdosproblems.com}

\begin{document}
\maketitle

\noindent\textbf{Status:} OPEN \quad \textbf{No prize}

\noindent\textbf{Tags:} number theory, iterated functions

\medskip

\noindent\textbf{Problem Statement:}

Let $\omega(n)$ count the number of distinct primes dividing $n$. Are there infinitely many $n$ such that, for all $m&lt;n$, we have $m+\omega(m) \leq n$? Can one show that there exists an $\epsilon&gt;0$ such that there are infinitely many $n$ where $m+\epsilon \omega(m)\leq n$ for all $m&#60;n$?




In \cite{Er79} Erd\H{o}s calls such an $n$ a 'barrier' for $\omega$. Some other natural number theoretic functions (such as $\phi$ and $\sigma$) have no barriers because they increase too rapidly. Erd\H{o}s believed that $\omega$ should have infinitely many barriers. In \cite{Er79d} he proves that $F(n)=\prod k_i$, where $n=\prod p_i^{k_i}$, has infinitely many barriers (in fact the set of barriers has positive density in the integers). Erd\H{o}s also believed that $\Omega$, the count of the number of prime factors with multiplicity), should have infinitely many barriers. Selfridge found the largest barrier for $\Omega$ which is $&lt;10^5$ is $99840$. In \cite{ErGr80} this problem is suggested as a way of showing that the iterated behaviour of $n\mapsto n+\omega(n)$ eventually settles into a single sequence, regardless of the starting value of $n$ (see also [412] and [414] ). Erd\H{o}s and Graham report it could be attacked by sieve methods, but 'at present these methods are not strong enough'. Lau \cite{La26} (see [679] ) has provided a positive answer to the second question, and proved a weaker version of the first: there exists a constant $C$ such that, for infinitely many $n$,\[m+\omega(m) \leq n\]for all $1\leq m\leq n-C$. The sequence of barriers for $\omega$ is A005236 in the OEIS. This is discussed in problem B8 of Guy's collection \cite{Gu04}. See also [647] and [679]




References

[Er79] Erd\H{o}s, Paul, Some unconventional problems in number theory . Math. Mag. (1979), 67-70.

[Er79d] Erd\H{o}s, P., Some unconventional problems in number theory . Acta Math. Acad. Sci. Hungar. (1979), 71-80.

[ErGr80] Erd\H{o}s, P. and Graham, R., Old and new problems and results in combinatorial number theory . Monographies de L'Enseignement Mathematique (1980).

[Gu04] Guy, Richard K., Unsolved problems in number theory . (2004), xviii+437.

[La26] C. F. Lau, On the number of prime factors of consecutive integers . arXiv:2604.15042 (2026).

\noindent\textbf{Related OEIS sequences:} A005236


\bigskip
\noindent\small{Source: \url{https://www.erdosproblems.com/413}}
\end{document}
